\subsection*{1. Laser absorption at maximum boundaries}
\phantomsection\label{bug:1}
Energy injected at \code{x_max}, \code{y_max} or \code{z_max}
could be recorded with the wrong sign. This could make the reported
absorption coefficient zero. The fix covers 1D, 2D and 3D laser code.
Tom Goffrey merged \pr{896} on 10 April 2026 as
\upcommit{0984b1d31}. It is in official \code{main} and in the
\href{https://github.com/epochpic/epoch/releases/tag/v4.20.1}{v4.20.1
release notes} (\issue{893}).

\subsection*{2. Linear Breit--Wheeler species selection in 2D}
\phantomsection\label{bug:2}
A one-line typo in the 2D linear Breit--Wheeler routine selected the
wrong species. \pr{892} merged \upcommit{ba04ae645} on
26 March 2026. The fix is in official \code{main} and listed in
v4.20.1 (\issue{891}). This is separate from the nonlinear
Breit--Wheeler items below.

\subsection*{3. Photon-probe particle-ID deadlock}
\phantomsection\label{bug:3}
\subsubsection*{When it happens}
With photon probes and \code{PARTICLE_ID} or \code{PARTICLE_ID4},
a rank with no photons could return before a call that all ranks
needed to make. Other ranks would wait indefinitely.

\subsubsection*{Where the fix is}
\pr{906} merged \upcommit{241d6fb9c} into the official
\code{4.20-devel} branch on 1 June 2026, in all three dimensions.
It is not in the official \code{main} version checked for this
document. A later 2D fork commit, \forkcommit{3e4734522}, also
changes this path but includes a separate ID fix (bug 4).

\subsubsection*{Test to keep}
Run with at least two MPI ranks, a photon probe, and one rank
without photons. Check that the run finishes and that the IDs
remain unique. One rank cannot expose this deadlock.

\subsection*{4. Particle IDs reused after restart or file loading}
\phantomsection\label{bug:4}
\textbf{Problem.} The 2D investigation found that the largest
assigned ID was not always carried through restart or particle
loading from a file. New particles could then reuse an old ID.

\textbf{Fix and status.} \forkcommit{3e4734522} saves and raises
the ID counter in the fork. The checked official \code{main}
does not contain those changes. The fork fix is 2D only.
A regression test should check IDs after particles leave the
domain, after a restart with a different MPI layout, and after
loading both supplied and automatic file IDs.

\subsection*{5. Higuera--Cary particle-push field scale}
\phantomsection\label{bug:5}
\subsubsection*{Problem}
The magnetic rotation in the Higuera--Cary particle push uses
an interpolated field. Its scale must include the same factor
used elsewhere in the push. In the checked official \code{main},
that factor is missing from this rotation.

\subsubsection*{Fix and test}
The fork fix \forkcommit{d86ceceab} changes the calculation in
1D, 2D and 3D. A separate branch based on official
\code{4.21-devel} contains \forkcommit{2684956bf}.
Neither change is in official \code{main}. A useful test is a
single particle in crossed fields, compared with an independent
trajectory, for several particle shapes. This document checks
the source; it does not rerun that test on the official build.

\subsection*{6. Trident pair creation with pairs disabled}
\phantomsection\label{bug:6}
\subsubsection*{Problem}
With \code{TRIDENT_PHOTONS} compiled and
\code{produce_pairs = F} in the deck, the electron loop could
still try to create a trident pair. If no pair species was set,
this could use species index \code{-1} and crash.

\subsubsection*{Fix and test}
The fork commit \forkcommit{1a88d4afa} checks
\code{produce_pairs} before changing the trident event counter
and before creating a pair, in all three dimensions.
The checked official \code{main} still has the unguarded path.
A test should run past a would-be event with pairs disabled and
confirm that the run completes without creating a pair.

\subsection*{7. Wrong photon species in nonlinear pair creation}
\phantomsection\label{bug:7}
The 2D photon loop in the checked official \code{main} passes
\code{photon_species} to pair creation instead of the species
currently being processed, \code{ispecies}. With more than one
photon species, this can use the wrong particle list and may
crash. The 2D fork fix \forkcommit{2b18ef6b4} passes
\code{ispecies}. A test needs at least two photon species
that can create nonlinear Breit--Wheeler pairs.

\subsection*{8. \code{identify} silently replaces charge or mass}
\phantomsection\label{bug:8}
In an EPOCH species block, \code{identify} can replace an
explicit \code{charge} or \code{mass} without warning. A deck
may then run with different particle properties than its
author expected. The fork commit \forkcommit{b908c8591}
adds a warning in all three dimensions. It does not change
which properties EPOCH uses. The warning is absent from the
checked official \code{main}. Tests should cover both input
orders and both matching and mismatching values.

\subsection*{9. Spin state lost on restart in the fork}
\phantomsection\label{bug:9}
This affects the fork's 2D spin feature; the feature is not in
the checked official \code{main}. The fork fix
\forkcommit{a66b394cb} restores stored spin state on restart.
A test should compare a continuous run with a restarted run.

\subsection*{10. Boosted pair-channel errors in the fork}
\phantomsection\label{bug:10}
This entry groups related fixes to pair-boost features absent
from the checked official \code{main}. They are separate code
changes: Breit--Wheeler boost handling
(\forkcommit{3e9f3121f}), muon gamma sampling
(\forkcommit{fbb9a9459}), Bethe--Heitler scaling
(\forkcommit{92552de69}, \forkcommit{a5716caf2}), and
additive event counters (\forkcommit{99ce3c685}).
Review and test each change separately before proposing it
for official EPOCH.

\subsection*{11. Half-cell offset in the fork's custom-laser reader}
\phantomsection\label{bug:11}
The fork's 2D file-based laser reader could sample a profile
half a cell away from the intended point. The correction is
\forkcommit{1b307b37d}. The checked official \code{main}
does not have this reader.

\subsection*{12. MPI deadlock in the fork's custom-laser reader}
\phantomsection\label{bug:12}
\subsubsection*{Problem and fix}
Ranks could take different paths while loading profile and
phase files, leaving some ranks waiting. The 2D fork fix is
\forkcommit{5a77c367b}. This concerns a reader absent
from the checked official \code{main}.

\subsubsection*{Test to keep}
Load both profile and phase files on multiple ranks, including
a rank with no requested local data. All ranks should finish
and the injected field should match the supplied data.

\subsection*{13. Two custom lasers sharing file data in the fork}
\phantomsection\label{bug:13}
The 2D fork reader could let one laser block use another
block's imported profile. \forkcommit{945fd060d} stores the
data with each laser block. The reader is absent from the
checked official \code{main}. A test should load two
different profiles on the same boundary and check both fields.

\subsection*{14. Stale 3D custom-laser profile after load balancing}
\phantomsection\label{bug:14}
In the fork's 3D reader, the loaded profile window needed to
follow changes to the domain split. \forkcommit{f572d035c}
reloads the relevant profile portion as the split moves.
The reader is absent from the checked official \code{main}.
A test should compare the injected field before and after
a load-balancing change.
